\documentclass[a4paper,oneside,14pt]{extreport}

% ==================================================
% Кодировки и языки (bmstu.cls)
% ==================================================
\usepackage{cmap}
\usepackage[OT2,T2C,T2B,T2A,LGR,T1]{fontenc}
\usepackage[utf8]{inputenc}
\usepackage[english,main=russian]{babel}
\usepackage{tempora}
\usepackage[scaled=0.95]{PTMono}

% ==================================================
% Геометрия страницы (bmstu.cls)
% ==================================================
\usepackage[
left=30mm,
right=15mm,
top=20mm,
bottom=20mm,
]{geometry}

% ==================================================
% Типографика и переносы
% ==================================================
\usepackage[expansion=false]{microtype}
\sloppy
\hyphenpenalty=50
\exhyphenpenalty=50

% ==================================================
% Межстрочный интервал
% ==================================================
\usepackage{setspace}
\onehalfspacing

% ==================================================
% Абзацы
% ==================================================
\usepackage{indentfirst}
\setlength{\parindent}{12.5mm}

% ==================================================
% Размеры шрифтов заголовков
% ==================================================
\makeatletter
\renewcommand\LARGE{\@setfontsize\LARGE{22pt}{20}}
\renewcommand\Large{\@setfontsize\Large{20pt}{20}}
\renewcommand\large{\@setfontsize\large{16pt}{20}}
\makeatother

% ==================================================
% Заголовки
% ==================================================
\usepackage{titlesec}

\titleformat{\chapter}[block]
{\hspace{\parindent}\large\bfseries}
{\thechapter}{0.5em}{\large\bfseries\raggedright}

\titleformat{name=\chapter,numberless}[block]
{\hspace{\parindent}}{}{0pt}{\large\bfseries\centering}

\titleformat{\section}[block]
{\hspace{\parindent}\large\bfseries}
{\thesection}{0.5em}{\large\bfseries\raggedright}

\titleformat{\subsection}[block]
{\hspace{\parindent}\large\bfseries}
{\thesubsection}{0.5em}{\large\bfseries\raggedright}

\titleformat{\subsubsection}[block]
{\hspace{\parindent}\large\bfseries}
{\thesubsubsection}{0.5em}{\large\bfseries\raggedright}

\titlespacing{\chapter}{12.5mm}{-22pt}{10pt}
\titlespacing{\section}{12.5mm}{10pt}{10pt}
\titlespacing{\subsection}{12.5mm}{10pt}{10pt}
\titlespacing{\subsubsection}{12.5mm}{10pt}{10pt}

% ==================================================
% Цвета
% ==================================================
\usepackage{xcolor,colortbl}

% ==================================================
% Таблицы и окружения (твои пакеты)
% ==================================================
\usepackage{tabularx}
\usepackage{tabulary}
\usepackage{threeparttable}
\usepackage{multirow}
\usepackage{booktabs}
\usepackage{pgfplots}
\pgfplotsset{compat=1.18}

% ==================================================
% Подписи
% ==================================================
\usepackage{caption}
\captionsetup{
	labelsep=endash,
	figurename=Рисунок,
	singlelinecheck=false,
}

\captionsetup[figure]{justification=centering}

\captionsetup[table]{
	justification=centering,
	singlelinecheck=false,
	position=above,
	skip=6pt
}

% ==================================================
% Повороты и альбомные страницы
% ==================================================
\usepackage{rotating}
\usepackage{lscape}
\usepackage{pdflscape}
\usepackage{afterpage}

% ==================================================
% Математика
% ==================================================
\usepackage{amsmath}
\usepackage{amssymb}

% ==================================================
% Библиография (bmstu.cls)
% ==================================================
\usepackage[
backend=biber,
style=gost-numeric,
language=auto,
autolang=other,
sorting=none,
]{biblatex}

\addbibresource{biblio.bib}

\usepackage{csquotes}

\DeclareFieldFormat{urldate}
{(Дата обращения:\addspace\ifnum\thefield{urlday}<10 0\fi\thefield{urlday}\adddot\ifnum\thefield{urlmonth}<10 0\fi\thefield{urlmonth}\adddot\thefield{urlyear})}
\DeclareFieldFormat[online]{title}{#1\addspace\mkbibbrackets{Электронный ресурс}}
\DeclareFieldFormat{url}{Режим доступа URL\addcolon\space\url{#1}}

% ==================================================
% Гиперссылки
% ==================================================
\usepackage[unicode,hidelinks]{hyperref}

% ==================================================
% Управляющие пакеты
% ==================================================
\usepackage{xifthen}
\usepackage{etoolbox}
\usepackage{totcount}
\usepackage{assoccnt}
\usepackage{placeins}
\usepackage{changepage}
\usepackage{pdfpages}
\usepackage{needspace}
%\usepackage{soul}
%\usepackage{svg}
%\usepackage{shellesc}

% ==================================================
% Содержание (bmstu-toc)
% ==================================================
\newcommand{\maketableofcontents}{
	\renewcommand\contentsname{СОДЕРЖАНИЕ}
	\tableofcontents
}

\setcounter{secnumdepth}{3}
\setcounter{tocdepth}{2}

% ==================================================
% Счётчики рисунков, таблиц, источников (bmstu-essay)
% ==================================================
\usepackage{lastpage}

\newcounter{totfigures}
\newcounter{tottables}
\providecommand\totfig{}
\providecommand\tottab{}

\makeatletter
\AtEndDocument{
	\addtocounter{totfigures}{\value{figure}}
	\addtocounter{tottables}{\value{table}}
	\immediate\write\@mainaux{
		\string\gdef\string\totfig{\number\value{totfigures}}
		\string\gdef\string\tottab{\number\value{tottables}}
	}
}
\makeatother

\pretocmd{\chapter}{\addtocounter{totfigures}{\value{figure}}\setcounter{figure}{0}}{}{}
\pretocmd{\chapter}{\addtocounter{tottables}{\value{table}}\setcounter{table}{0}}{}{}
\pretocmd{\section}{\Needspace{8\baselineskip}}{}{}
\pretocmd{\subsection}{\Needspace{7\baselineskip}}{}{}
\pretocmd{\subsubsection}{\Needspace{6\baselineskip}}{}{}

% ==================================================
% Подсчёт источников
% ==================================================
\newcounter{totbibentries}
\newcommand*{\listcounted}{}

\makeatletter
\AtDataInput{
	\xifinlist{\abx@field@entrykey}\listcounted
	{}
	{\stepcounter{totbibentries}
		\listxadd\listcounted{\abx@field@entrykey}}
}
\makeatother

% ==================================================
% Приложения (bmstu-appendix)
% ==================================================
\usepackage[titletoc,title]{appendix}
\usepackage{titlesec}

\AtBeginDocument{\renewcommand{\appendixname}{ПРИЛОЖЕНИЕ}}

\let\oldappendices\appendices
\let\oldendappendices\endappendices

% Новое окружение appendices с опциональным аргументом для буквы
\usepackage{xparse}
\RenewDocumentEnvironment{appendices}{o} % o = optional argument
{
	\titleformat{\chapter}{\large\bfseries}
	{\appendixname~\IfValueTF{#1}{#1}{\thechapter}}{0pt}
	{\centering\large\bfseries\\}
	\oldappendices
	% если передали букву, используем её для визуала и PDF bookmarks
	\IfValueTF{#1}{
		\renewcommand{\thechapter}{#1}
	}{
		\renewcommand{\thechapter}{\Asbuk{chapter}} % кириллическая нумерация по умолчанию
	}
}
{
	\oldendappendices
}
% ==================================================
% Определения и сокращения (твоя версия)
% ==================================================
\usepackage{enumitem}

\newcommand{\mydefinition}[2]{
	\item \noindent #1 --- #2
}

\newenvironment{mydefinitions}{
	\chapter*{ОПРЕДЕЛЕНИЯ}
	\addcontentsline{toc}{chapter}{ОПРЕДЕЛЕНИЯ}
	\begin{description}[leftmargin=0pt]
	}{
	\end{description}
}

\newenvironment{myabbreviations}{
	\chapter*{ОБОЗНАЧЕНИЯ И СОКРАЩЕНИЯ}
	\addcontentsline{toc}{chapter}{ОБОЗНАЧЕНИЯ И СОКРАЩЕНИЯ}
	\begin{description}[leftmargin=0pt]
	}{
	\end{description}
}

% ==================================================
% TOC: оформление глав (твой патч)
% ==================================================
\makeatletter
\renewcommand*{\l@chapter}[2]{
	\ifnum \c@tocdepth>\m@ne
	\addpenalty{-\@highpenalty}
	\vskip 1em \@plus.2em
	\@dottedtocline{0}{0pt}{1.5em}{\bfseries #1}{\bfseries #2}
	\fi
}
\makeatother


\newenvironment{essay}[1]
{
	\clearpage
	\chapter*{РЕФЕРАТ}
	\addcontentsline{toc}{chapter}{РЕФЕРАТ}
	
	% --- строка с подсчётами (безопасно, если приложений ещё нет)
	Расчётно-пояснительная записка \begin{NoHyper}\pageref{LastPage}\end{NoHyper}~с., 
	\totfig~рис., \tottab~табл., \arabic{totbibentries}~источн., 2~прил.
	
	\noindent\MakeUppercase{#1} \par
}
{
	%\vspace{1em}
	\clearpage
}

% ==================================================
% Пользовательские команды
% ==================================================
\newcommand*{\rom}[1]{\uppercase\expandafter{\romannumeral #1\relax}}

\newenvironment{formulaexplanation}
{
	\par\noindent где
	\begin{itemize}[label=---, itemsep=0pt, parsep=0pt, topsep=0pt]
}
{
	\end{itemize}
}
